
\documentclass[10pt,a4paper,english]{report}
\usepackage{amssymb}




%%%%%%%%%%%%%%%%%%%%%%%%%%%%%%%%%%%%%%%%%%%%%%%%%%%%%%%%%%%%%%%%%%%%%%%%%%%%%%%%%%%%%%%%%%%%%%%%%%%%
%\usepackage{makeidx}
%\usepackage[T1]{fontenc}
%%\usepackage[latin9]{inputenc}
\usepackage{amsmath}
\usepackage{graphicx}
\usepackage{graphics}
\usepackage[dvips]{epsfig}
\usepackage{esint}
\usepackage[authoryear]{natbib}
\usepackage{rotating}
\usepackage{relsize}
\usepackage{Sweave}
\usepackage{xcolor}
\usepackage{babel}
\usepackage{hyperref}



\begin{document}
\title{Computer Exercise 4: Linear Mixed Models with Correlated Random Effects}
\author{Lars Ronnegard}
\maketitle
\section*{Introduction}
The first four exercises focusses on the animal model using two data sets with pedigree information. In the last two exercises your task is to explore semiparametric smoothing. \newline
\newline
The first example data set we will use to explore the animal model is the \textbf{l51} data set in the gap package. There are observations on a quantitative trait from 51 related individuals. We will use the hglm package to estimate the additive genetic variance and the residual variance. \newline
\newline
An animal model is a linear mixed model with correlated random effects where $y = X\beta + Za + e$ and $V(a)=\mathbf{A}\sigma^2_a$. $\mathbf{A}$ is the additive relationship matrix. 
\newline
\newline
The pedigree in the \textbf{l51} data set can be plotted using the following code.

\begin{Schunk}
\begin{Sinput}
library(gap)
data(l51, package="gap")
library(kinship2)
ped <- with(l51, pedigree(id,fid,mid,sex))
plot(ped)
\end{Sinput}
\end{Schunk}

\begin{figure}[h]
\center \includegraphics[width=1.0\textwidth]{L51_pedigree_plot}
\caption{Pedigree for the \textbf{l51} data.}
\label{fig:L51_pedigree_plot}
\end{figure}
\newpage 

The following code shows how the additive relationship matrix $\mathbf{A}$ iin an animal model can be computed using the pedigree information \texttt{l51} and the kin.morgan() function in the \textbf{gap} package.

\begin{Schunk}
\begin{Sinput}
km <- kin.morgan(l51)
A <- km$kin.matrix*2
image(A)
\end{Sinput}
\end{Schunk}

\begin{figure}[h]
\center \includegraphics[width=0.8\textwidth]{L51_A_image}
\caption{Image plot of the additive relationship matrix for the \textbf{l51} pedigree data.}
\label{fig:L51_A_image}
\end{figure}

\newpage 

The following code shows how an animal model can be fitted using the hglm package. If you get stuck, have a look at the YouTube tutorials by Xia Shen at  \url{https://www.youtube.com/user/UppsalaQTL/featured}.

\begin{Schunk}
\begin{Sinput}
########## hglm analysis ##########
library(hglm)
obs.qt <- !is.na(l51$qt) #Remove missing data
L <- t( chol( A ) )[ obs.qt, ] #Compute 
X <- model.matrix( qt ~ 1, data = l51)
y <- l51$qt[ obs.qt ]
mod1 <- hglm(y = y,  X = X, Z = L)
summary(mod1)

\end{Sinput}
\end{Schunk}

\section*{Exercise 1}
In a linear mixed model with independent random effects the variance covariance matrix for the response $y$ is: $V(y)=\mathbf{Z}\mathbf{Z}^T\sigma^2_a+\mathbf{I}\sigma^2_a$. \newline
In an animal model we have: $V(y)=\mathbf{A}\sigma^2_a+\mathbf{I}\sigma^2_a$. \newline
The two models are equivalent if $\mathbf{Z}=\mathbf{L}$ and $\mathbf{A}=\mathbf{L}\mathbf{L}^T$. \newline
Check that  $\mathbf{A}=\mathbf{L}\mathbf{L}^T$ in our example.


\section*{Exercise 2}
Identify the variance component estimates from the hglm output. Compute the estimated heritability.

\section*{Exercise 3}
Compute the estimated breeding values for all individuals in the pedigree using the hglm results.

\section*{Exercise 4 (If you get time)}
Install the MCMCglmm package.
Combine the two data sets BTped and BTdata in the package using the code below, and fit an animal model for tarsus length as response and sex as explanatory variable.
\newline
Use the hglm package, but first compute the design matrix X, the incidence matrix Z and its Cholesky factorization L.
\newline
Make an image plot of the additive relationship matrix.
\newline
Make diagnostic plots.
\newline
Estimate the heritability for tarsus length.

\begin{Schunk}
\begin{Sinput}
data(BTped, package="MCMCglmm")
head(BTped)
tail(BTped)
data(BTdata, package="MCMCglmm")
head(BTdata)
summary(BTdata)

invA <- inverseA(BTped)$Ainv
A <- solve(invA)
L <- t( chol(as.matrix(A)) )
obs.tarsus = 213:1040
ord.data <- merge(BTdata, BTped[obs.tarsus,], sort=FALSE)
Z0 <- diag(1040)[obs.tarsus,] 
Z <- Z0 %*% L
y <- ord.data$tarsus
\end{Sinput}
\end{Schunk}

\section*{Exercise 5}
In this exercise you will explore how fossil age depends on the proportion of strontium in the fossil. First plot the data using the following code. 
\begin{Schunk}
\begin{Sinput}
library(SemiPar)
data(fossil)
?fossil
plot(x=fossil$age, y=fossil$strontium.ratio)
\end{Sinput}
\end{Schunk}


To be able to perform semiparametric smoothing with the hglm package, we need a function to construct the spatial Matern covariance matrix, which depends on the distance between x values, a range parameter and a dimension parameter.
\begin{Schunk}
\begin{Sinput}
##Function to construct the Matern covariance matrix
get.matern <-function(d, l, nu) {
  #Distance matrix d
  #Range l
  #Dimension parameter nu
  if (nu == 0.5) c <- exp(-d*l^(-1))
  if (nu == 1.5) c <- (1+sqrt(3)*d*l^(-1))*exp(-sqrt(3)*d*l^(-1))
  if (nu == 2.5) c <- (1+(sqrt(5)*d)*l^(-1)+(5*d^2)*3^(-1)*l^(-2))*exp(-sqrt(5)*d*l^(-1))
  return(c)
}
##
\end{Sinput}
\end{Schunk}

We also need a function to make sure that the covariance matrix is positive definite and thereby has good numerical properties.
\begin{Schunk}
\begin{Sinput}
#see Ranjan et al (2011) Technometrics, 53:4, 366-378
bend <- function(C, OK.cond = 5) {
  eigVal <- eigen(C)$values
  #plot(-log10(eigVal))
  cond0 <- log10(max(eigVal)/min(eigVal)) 
  delta = max(eigVal)*(10^cond0-10^OK.cond)/(10^cond0*(10^OK.cond-1))
  C1 <- C
  diag(C1) = diag(C) + delta
  return(C1)
}
\end{Sinput}
\end{Schunk}

Before performing the analysis  the data is sorted according to age and the response  is scaled to have a variance of 1 just to simplify plotting at the end.
\begin{Schunk}
\begin{Sinput}
x <- as.numeric(sort(fossil$age))
y <- scale(fossil$strontium.ratio[sort(fossil$age, index.return=TRUE)$ix])
plot(x,y, xlab="Age", ylab="Scaled Strontium Ratio")
\end{Sinput}
\end{Schunk}

The following code fits the model and plots the smoothed function estimated by the hglm package.

\begin{Schunk}
\begin{Sinput}
distMat = as.matrix(dist(x))
rho = 9.1 #A good first value to try is sd(x)
nu = 1.5 #Can be either 0.5, 1.5 or 2.5
C1 <- get.matern(distMat, rho, nu)
C1 <- bend(C1) 
image(C1)
Z <- t(chol(C1))
N <- length(y)
X <- matrix(1, N, 1)
library(hglm)
mod1 <- hglm(y=y, X=X, Z=Z)
plot(x,y, xlab="Age", ylab="Scaled Strontium Ratio")
points(x , mod1$fv, col="red", type="l")
\end{Sinput}
\end{Schunk}
Try different values of \texttt{rho} and \texttt{nu}. Which combination of values seem most reasonable? \newline
What does the image plot in Figure 3 show?

\begin{figure}[h]
\center \includegraphics[width=0.8\textwidth]{FossilC1_plot}
\caption{Image plot of the spatial covariance matrix for the fossil data.}
\label{fig:spatial}
\end{figure}

\section*{Exercise 6 (If you get time)}
Which function is hiding behind the following data??? 
\begin{Schunk}
\begin{Sinput}
x <- 1:100
y <- c(8.0976, 9.5571, -0.3118, -4.4153, 8.1990, -20.4927, 8.4122, -2.2327, 
   -2.3094, -9.6534, -0.3404, 21.8501, 5.2310, 6.9091, -7.4827, 10.1867, -8.9181, 
   -2.6918, 13.4157, 4.2929, 9.5510, 16.8877, -5.9968, -15.7815, -16.1770, 
   20.3098, -4.1267, 7.5704, 7.4506, -0.7019, 9.4705, 23.8058, 22.1670, 17.7782, 
   3.6004, 5.9982, -2.2814, -15.7442, 18.7201, 1.2573, 12.1998, -22.5059,
   -9.4178, 10.1621, 9.3120, 15.1390, -12.6097, 6.4375, 6.5609, -11.3986, 
   -4.5070, 20.2074, 3.3301, 5.5128, -6.4163, -1.2048, -0.4397, 0.0819, 14.7244, 
   -20.3520, 8.5935, -0.2709, 16.7784, 20.9254, -7.0288, -27.0844, -9.0063, 
   -2.2829, -4.1556, -7.8179, 15.5696, 15.1409, 11.2536, 29.5960, 6.5936, 
   -11.5797, -20.0413, -30.6142, -7.0520, 15.7339, 29.2483, 40.3961, 25.0073, 
   31.3611, -2.2912, -14.0007, -31.6998, -33.1300, -7.5069, 19.4468, 26.3393, 
   35.3989, 48.1140, 8.4584, -23.8512, -33.7772, -50.6525, -20.2156, -15.2219, 
   15.0303)
par(mfrow=c(1,1))
plot(x=x, y=y)
\end{Sinput}
\end{Schunk}

\end{document}